\documentclass[10pt,a4paper]{amsart}
\usepackage[margin=19mm]{geometry}
\usepackage{amsmath,amssymb,mathtools}
\usepackage[scr=rsfs,cal=euler]{mathalfa}
\usepackage{xfrac}
\usepackage[unicode,hidelinks,bookmarks=false]{hyperref}
\newcommand{\ie}{i.e.\ }
\newcommand{\cf}{cf.\ }
\newcommand{\resp}{respectively\ }
\newcommand{\Subsection}{\subsection}
\providecommand{\nobreakdash}{\nobreak\hbox{-}}
\setlength{\emergencystretch}{2em}
\multlinegap0pt
\usepackage[shortlabels]{enumitem}
\usepackage{amssymb}
\usepackage{mathtools}
\usepackage{algorithm}\usepackage{algpseudocode}
\usepackage{float}
\usepackage{tikz}
\newtheorem{lemma}{Lemma}
\newtheorem{definition}{Definition}
\newtheorem{theorem}{Theorem}
\newcommand{\Cn}{\text{Cn}}
\newcommand{\Dd}{\text{Dd}}
\title{The Derivation Penalty in Premise-Erasure Caching: Capacity, Strong Converse, and Dispersion Dichotomy}
\author{Jianfeng Xu}
\date{Source: arXiv:2603.00930v1 (CC BY 4.0)}
\begin{document}
\maketitle

\noindent\emph{Source and licence.} The Derivation Penalty in Premise-Erasure Caching: Capacity, Strong Converse, and Dispersion Dichotomy. Version arXiv:2603.00930v1 (CC BY 4.0), distributed by arXiv under the Creative Commons Attribution 4.0 International licence, http://creativecommons.org/licenses/by/4.0/, as recorded in the arXiv licence metadata for this exact version and shown on the abstract page https://arxiv.org/abs/2603.00930v1. Full licence notice: \texttt{licenses/arxiv-2603.00930-CC-BY-4.0.txt}.\par\smallskip

\noindent\textit{Editorial scope.} Counted unit: the article's theorem thm:loss-comp-duality (source0.tex lines 1353--1377). The article's complete original proof of it is delivered with it (source0.tex lines 1379--1412, one \texttt{proof} environment of 34 lines). No mathematics is written, repaired or re-derived: the statement and the proof are the article's own lines, in the article's own notation, and no retained line is substituted or re-flowed.  Retained uncounted as the source-local context the proof needs: lines 1245--1249; lines 1343--1351.  Nothing was omitted from any retained span, so the package carries no omission record.  No retained line contains a citation, so the wrapper carries no bibliography and no bibliography entry.  No retained line contains a figure environment, an image-inclusion command or drawing code, so no image is delivered and the package declares no asset dependency.  Editorial formatting only, in addition: an AMS \texttt{amsart} preamble that restates the article's own declarations and macros used by the retained lines, namely the environments \texttt{lemma}, \texttt{definition}, \texttt{theorem} on separate counters, together with the standard packages those lines need.  The retained environments print in this document as Lemma 1, Definition 1, Theorem 1 in order of appearance, whereas the article prints the same items with the numbering of its own full document.  The complete native source of the article as distributed in the arXiv e-print for this exact version is bundled unchanged as \texttt{sources/195527/source0.tex}.
\begin{lemma}[Monotonicity in the premise base]
\label{lem:premise-monotonicity-sec3}
If $B_1\subseteq B_2$, then $\Dd(s\mid B_2)\le\Dd(s\mid B_1)$
for every $s\in\mathbb S_O$.
\end{lemma}

\begin{definition}[Global reconstruction depth]
\label{def:rec-depth-preserved}
Let $B_0^-\subseteq B_0$ and $B_\cap:=B_0\setminus B_0^-$.
The \emph{reconstruction depth} is
$d_{\mathrm{rec}}:=\max_{b\in B_0^-}\Dd(b\mid B_\cap)$,
with $d_{\mathrm{rec}}:=0$ if $B_0^-=\varnothing$ and
$d_{\mathrm{rec}}:=\infty$ if some lost premise is not
derivable from $B_\cap$.
\end{definition}

\begin{theorem}[Loss--computation duality]
\label{thm:loss-comp-duality}
Let $B_0^-\subseteq B_0$, $B_\cap=B_0\setminus B_0^-$,
$q\in\Cn(B_\cap)$.

\begin{enumerate}[label=\textup{(\Roman*)}]
\item \emph{Depth shift.}
\begin{equation}\label{eq:depth-shift}
  \Dd(q\mid B_0)
  \;\le\;
  \Dd(q\mid B_\cap)
  \;\le\;
  \Dd(q\mid B_0)+d_{\mathrm{rec}}.
\end{equation}

\item If $B_0^-\cap V_0(G(q,B_0))=\varnothing$
\textup{(no lost fact appears in the dependency set)}, then
$\Dd(q\mid B_\cap)=\Dd(q\mid B_0)$.

\item Under combined noise
$\tilde B_0=(B_0\setminus B_0^-)\cup B_0^+$:
$\Dd(q\mid\tilde B_0)\le\Dd(q\mid B_\cap)
 \le\Dd(q\mid B_0)+d_{\mathrm{rec}}$.
\end{enumerate}
\end{theorem}

\begin{proof}
(I)~The lower bound is
Lemma~\ref{lem:premise-monotonicity-sec3} ($B_\cap\subseteq B_0$).
For the upper bound we prove
$\Dd(v\mid B_\cap)\le\Dd(v\mid B_0)+d_{\mathrm{rec}}$
for every $v$ in the backward unfolding of~$q$,
by well-founded induction on $\Dd(v\mid B_0)$.

\emph{Base.}\;$v\in B_0$.
If $v\in B_\cap$, then
$\Dd(v\mid B_\cap)=0\le 0+d_{\mathrm{rec}}$.
If $v\in B_0^-$, then
$\Dd(v\mid B_\cap)\le d_{\mathrm{rec}}
 =0+d_{\mathrm{rec}}=\Dd(v\mid B_0)+d_{\mathrm{rec}}$
by Definition~\ref{def:rec-depth-preserved}.

\emph{Step.}\;$v\notin B_0$.
Each $v'\in P_O(v)$ satisfies
$\Dd(v'\mid B_0)<\Dd(v\mid B_0)$; by the induction
hypothesis,
$\Dd(v'\mid B_\cap)\le\Dd(v'\mid B_0)+d_{\mathrm{rec}}$.
Hence
$\Dd(v\mid B_\cap)
 =1+\max_{v'}\Dd(v'\mid B_\cap)
 \le\Dd(v\mid B_0)+d_{\mathrm{rec}}$.
Specializing to $v:=q$ yields~\eqref{eq:depth-shift}.

(II)~The derivation of $q$ from $B_0$ uses only premises in
$B_\cap$; equality follows from
Lemma~\ref{lem:premise-monotonicity-sec3}.

(III)~$B_\cap\subseteq\tilde B_0$ gives the first inequality by
Lemma~\ref{lem:premise-monotonicity-sec3}; the second is from~(I).
\end{proof}

\end{document}
